\chapter*{Abstract} \label{abtract}

\addcontentsline{toc}{chapter}{Abstract} % add to index

Modern cloud infrastructures increasingly host specialized, single-application workloads, but these systems still run on a general-purpose Linux kernel configured through broad default \textit{sysctl} settings.
While these default configurations ensure stability and portability, they are frequently suboptimal for high-performance services.
Because the kernel configuration space is massive and complex, manual tuning is largely impractical.
This thesis addresses that gap by proposing an application-agnostic methodology to automatically identify which kernel knobs will actually affect a target workload before the optimization process begins.

To achieve this, the core contribution is a static analysis pipeline that combines two complementary views of kernel behaviour.
First, it performs a top-down reconstruction of the execution paths triggered by a system call.
Second, it uses a bottom-up analysis to map where specific \textit{sysctl} knobs are evaluated within the kernel code.
Intersecting these two views filters out irrelevant parameters, transforming a highly complex, high-dimensional search space into a drastically reduced subset of knobs suitable for automated optimization.

The methodology is implemented and evaluated on web-server workloads (Apache HTTP Server, NGINX, and HAProxy).
To quantify performance, we measured application throughput (Requests Per Second) under heavy concurrency, comparing the optimized configurations against the default, un-tuned Linux kernel baseline.
Throughput was selected as the primary metric because it aggressively stresses the kernel network stack and connection management subsystems.
By coupling the targeted, reduced search space with an existing autotuner, we significantly decreased the overall tuning time required for the algorithm to reach convergence.
Furthermore, the system successfully discovered configurations that yielded substantial throughput gains: up to 23\% for Apache and approximately 15\% for NGINX and HAProxy.

Overall, the experimental results confirm that static analysis of kernel paths can effectively and efficiently guide automated tuning toward meaningful performance improvements.
Ultimately, this work bridges the gap between static operating system analysis and automated tuning, providing a practical foundation for scalable, autonomous kernel optimization in modern virtualized environments.

\newpage
